% =====================================================================
%  III. NISQ-Aware QSVM Framework  (phan dau; theory_revision.tex noi tiep)
%
%  Muc nay chi dung khung: dat bai toan, ta pipeline, va neu ro gia thiet
%  cost model. Phan ly thuyet (F1-F3, Lemma 1, luat ba giai doan, erratum)
%  nam trong theory_revision.tex ngay sau day.
% =====================================================================

\subsection{Feature map, kernel and cost}

The \ZZ{}FeatureMap on $n$ qubits with $r$ repetitions prepares
\begin{equation}
\label{eq:featuremap}
  |\phi(\bm{x})\rangle \;=\;
  \bigl[\,U_{\Phi}(\bm{x})\,H^{\otimes n}\,\bigr]^{r}\,|0\rangle^{\otimes n},
\end{equation}
where $H$ is the Hadamard gate and $U_{\Phi}$ is diagonal in the computational
basis,
\begin{equation}
\label{eq:phase}
  U_{\Phi}(\bm{x}) \;=\;
  \exp\!\Bigl(\mathrm{i}\!\sum_{i=1}^{n}\! x_i Z_i
  \;+\; \mathrm{i}\!\!\sum_{1\le i<j\le n}\!\!
  \tfrac{1}{2}\varphi_{ij}(\bm{x})\,Z_iZ_j\Bigr),
  \quad
  \varphi_{ij}(\bm{x}) = 2(\pi-x_i)(\pi-x_j).
\end{equation}
The induced kernel is the squared fidelity
\begin{equation}
\label{eq:kernel}
  K^{\mathrm{Q}}(\bm{x},\bm{x}^{\prime}) \;=\;
  \bigl|\langle\phi(\bm{x}^{\prime})|\phi(\bm{x})\rangle\bigr|^{2},
\end{equation}
and the circuit cost we charge for it is
\begin{equation}
\label{eq:cost}
  Q(n) \;=\; \frac{Q_{\mathrm{raw}}(n)}{Q_{\mathrm{raw}}(n_{\max})},
  \qquad Q_{\mathrm{raw}}(n) = 10n^{2}-8n .
\end{equation}
Removing the pairwise term of Eq.~\eqref{eq:phase} while leaving everything else
unchanged gives the \Zmap{}FeatureMap, which is the control used throughout the
entanglement ablation.

\subsection{Problem statement}

Let $\mathcal{D}=\{(\bm{x}_i,y_i)\}_{i=1}^{N}$ be a labelled sample of network
records with $\bm{x}_i\in\mathbb{R}^{d}$ and $y_i$ one of five classes, and let
$K^{\mathrm{Q}}_{n}$ denote the kernel of Eq.~\eqref{eq:kernel} on an $n$-qubit
map. A quantum kernel is not free: evaluating $K^{\mathrm{Q}}_{n}$ over a
training set of size $N$ requires $\Theta(N^2)$ circuit evaluations, each
costing $Q(n)$ two-qubit gates on hardware whose error rate is dominated by
exactly those gates. The question this paper answers is therefore not whether
$K^{\mathrm{Q}}_{n}$ can be made to win somewhere, but where it is worth this
cost.

\begin{assumption}[NISQ cost model]
\label{ass:cost}
Circuit cost is charged by width alone, through the CNOT-weighted count of
Eq.~\eqref{eq:cost}, normalised by its value at the largest candidate so
that $Q(n)\in(0,1]$. The quantity it tracks is the
entangling-gate budget: an $r$-repetition \ZZ{}FeatureMap with full entanglement
places $2r\binom{n}{2}$ two-qubit gates, which at the operating point $n=4$,
$r=2$ is $24$ and at $n=8$ is $112$. The $\Theta(N^{2})$ kernel evaluations are
excluded from $Q$ because they scale with $N$ rather than $n$ and so do not
affect the choice of width; they are, however, the dominant cost of deploying
the method at all, which is why Sec.~\ref{sec:regimemap} reports the $N$ at
which the kernel becomes competitive and not merely that it does.
\end{assumption}

Assumption~\ref{ass:cost} is what makes the selection rule of
Definition~\ref{def:three-stage} a statement about hardware rather than about
taste, and it is also its main limitation: a device whose dominant error source
is idle-time decoherence rather than entangling-gate infidelity would call for a
different $Q$.

\begin{problem}[Regime identification]
\label{prob:regime}
Given a family of operating conditions --- training-set size, class prior,
attack composition, test distribution, and embedding width --- determine, for
each condition, whether $K^{\mathrm{Q}}_{n^\ast}$ is preferable to a classical
baseline, at a stated level of statistical confidence, and report the conditions
in which it is not.
\end{problem}

Problem~\ref{prob:regime} is deliberately weaker than the claim the submitted
version made. It admits negative answers as results, and most of the answers we
obtain are negative or inconclusive.

\subsection{Pipeline}

Fig.~\ref{fig:pipeline} shows the full path from raw records to a decision.
Categorical fields are one-hot encoded to $d=122$ dimensions; an ANOVA
$F$-test retains $K$ features; PCA projects to $n$ components; and a min--max
scaler maps each component into $[0,\pi]$, the domain on which the feature map's
phases are defined. The resulting $\bm{z}\in[0,\pi]^{n}$ enters the
\ZZ{}FeatureMap of Fig.~\ref{fig:circuit}, and classification is a soft-margin
SVC on the precomputed Gram matrix.

Two properties of this pipeline carry the experimental design.

\emph{One component varies at a time.} The entanglement-free control shown as
the upper branch of Fig.~\ref{fig:pipeline} differs from the main path in
exactly one block: the \Zmap{}FeatureMap applies the same Hadamard and
single-qubit phase layers and omits the pairwise term. It has the same qubit
count, the same repetitions, the same classical front-end, and the same
classifier, so a difference between the two branches is attributable to
entanglement and to nothing else. Fig.~\ref{fig:contributions} extends this
discipline to the other three contributions: C1 fixes the representation before
any downstream metric is inspected, C3 varies only the test distribution, and C4
varies only $N$.

\emph{The front-end is refit wherever the classifier is refit.} In the primary
arm, feature selection, projection and scaling are re-estimated at every
training-set size on that size's training fold alone. This costs comparability
across sizes --- the representation at $N=100$ is not the representation at
$N=10^{4}$ --- and buys the guarantee that no representation is ever informed by
more data than the classifier it feeds. We regard that trade as mandatory in a
study whose headline concerns how behaviour changes with $N$: freezing a
front-end fit on the largest subset would leak precisely the quantity under
investigation. Sec.~\ref{sec:res-c4} reports both arms.

\subsection{Kernel evaluation}

The kernel of Eq.~\eqref{eq:kernel} is evaluated exactly rather than by
sampling. Because the \ZZ{}FeatureMap is
diagonal in the computational basis after each Hadamard layer, $|\phi(\bm{x})
\rangle$ has a closed form, and the Gram matrix follows from dense linear
algebra. Finite-shot estimation and backend-derived noise are then applied as
\emph{separate} conditions (Sec.~\ref{sec:res-c2}) rather than being folded into
the main results, which keeps sampling error from being confused with the
effects under study.
